% Ampliación aditiva 2026-09-27. Propietarios y verificadores: recibo editorial contiguo.
\section{Composition d’action, d’aire et de lecture thermique marquée}
\label{sec:ii-termica-marcada}

La construction précédente détermine un porteur gradué, sa mémoire
énergétique, les courants entre secteurs et une référence spectrale marquée.
Cette section établit comment ces données agissent sur les lecteurs d’énergie
et d’information de l’article. Le développement comporte deux niveaux démonstratifs :
la composition des sections d’action, d’aire et d’horizon déjà construites,
et une réalisation thermométrique explicite au moyen de deux fonctionnelles du même
registre étendu. La seconde opération fixe un coefficient thermique dans
la règle de lecture énoncée et démontrée ici.

Les antécédents conservent la séquence APP--TRIT--TPK : les feuillets de somme
et de produit retiennent le résidu et le quotient ; TRIT conserve le régime et l’orientation ;
TPK transporte le chemin, la frontière, la retenue et la mémoire. La connexion nonadique ramène
la phase tout en prolongeant l’état enrichi. Les coordonnées
$\pi_{\HMT},\varphi_{\HMT},e_{\HMT},\alpha_{\HMT}$, l’action
orientée et l’incidence hexade--octade sont des sorties antérieures de cette même
structure discrète conjointe du continu. Dans les équations suivantes,
$\pi,\alpha$ et les autres lecteurs scalaires conservent cette provenance.

\subsection{Énergie de retour et résolution aréale}
\label{subsec:ii-termica-accion-area}

Fixons une orientation $a$ et la réalisation radiale de l’article. Nous écrivons
\begin{equation}
 L_\alpha=\frac{5759}{23040}\alpha^{16}L_* ,\qquad
 T_{108}=108t_0=\frac{2\pi L_\alpha}{c},\qquad
 \epsilon_a=\frac{h_a}{T_{108}}=\frac{\hbar_a c}{L_\alpha}.
 \label{eq:ii-ter27-energia}
\end{equation}
Le symbole $L_\alpha$ désigne une longueur ; le degré énergétique $L$ du
porteur précédent est un opérateur adimensionnel distinct. La fonctionnelle
orientée et le caractère angulaire précédemment construits sont
\begin{align}
 \gamma&=\frac{\pi AC^*}{180}
 +12\bigl[S_{90}(q_-)-S_{90}(q_+)\bigr]
 -\bigl[S_{120}(q_-)-S_{120}(q_+)\bigr],\nonumber\\
 S_m(q)&=\frac{q^m}{1-q^{3m}},\qquad
 a_0=\frac{1+2\cos(10^\circ)}3\frac{C^*}{6}\frac{\pi}{180}.
 \label{eq:ii-ter27-barbero}
\end{align}
Dans son secteur positif, le quantum d’aire est
$q_A=\gamma a_0L_\alpha^2$. Définissons l’énergie par unité d’aire
$\mathcal T_{I/A}=\epsilon_a/q_A$.

\begin{proposition}[Composition aréale, constitutive et gravitationnelle]
\label{prop:ii-ter27-area}
Avec le coefficient radial $G_a\hbar_a=c^3L_\alpha^2$, on a
\begin{align}
 \mathcal T_{I/A}&=\frac{c^4}{\gamma a_0G_aL_\alpha},&
 \frac{\mathcal T_{I/A}q_A}{\ln3}
 &=\frac{h_a}{108t_0\ln3},\label{eq:ii-ter27-area-energy}\\
 G_a\,k_B\Theta_{{\rm clk},a}\ln3&=c^4L_\alpha.&
 \label{eq:ii-ter27-grav-clock}
\end{align}
Si $\mathfrak e_a$ désigne la charge de la section constitutive
$\mathfrak e_a^2=2\alpha h_a/Z_{\rm vac}$, la même énergie par nat est
\begin{equation}
 \frac{\epsilon_a}{\ln3}
 =\frac{Z_{\rm vac}\mathfrak e_a^2}{216\alpha t_0\ln3}.
 \label{eq:ii-ter27-constitutive}
\end{equation}
\end{proposition}
\begin{proof}
Substituer $q_A$ et \eqref{eq:ii-ter27-energia} donne
$\mathcal T_{I/A}=\hbar_ac/(\gamma a_0L_\alpha^3)$.
L’identité radiale transforme cette expression en la première de
\eqref{eq:ii-ter27-area-energy} ; multiplier par $q_A$ prouve la seconde.
La relation thermique déjà établie est
$k_B\Theta_{{\rm clk},a}\ln3=\epsilon_a$.
Son produit par $G_a$ annule $\hbar_a$ et prouve
\eqref{eq:ii-ter27-grav-clock}. Enfin,
$h_a=Z_{\rm vac}\mathfrak e_a^2/(2\alpha)$ dans
\eqref{eq:ii-ter27-energia} produit \eqref{eq:ii-ter27-constitutive}.
\end{proof}

Barbero détermine la résolution aréale qui intervient dans la tension. Le
produit des deux lectures retrouve l’énergie de retour. Cette annulation
explique la compatibilité entre les formulations géométrique, constitutive
et thermique : toutes trois transportent la même section d’action.

\subsection{Information de provenance et réponse d’horizon}
\label{subsec:ii-termica-horizonte}

Sur les trente-six facteurs tritiques de bord, nous conservons
$N=N_{90}+N_{120}$ et $D=90N_{90}+120N_{120}$. L’aire est
$\widehat{\mathcal A}=q_AN$, tandis que l’état sectoriel et son
hamiltonien modulaire sont
\begin{equation}
 \rho_A=Z_A^{-1}e^{-\Delta D},\qquad
 K_A=-\ln\rho_A=(\ln Z_A)I+\Delta D.
 \label{eq:ii-ter27-modular}
\end{equation}
Ici, $\Delta$ est fini et $Z_A$ est la normalisation connue. L’incidence
permet de retrouver
\[
 N_{90}=4\widehat{\mathcal A}/q_A-D/30,\qquad
 N_{120}=D/30-3\widehat{\mathcal A}/q_A.
\]
Pour obtenir $D$ à partir de $K_A$, on utilise en outre $\Delta\ne0$.
Les configurations $(1,0)$ et $(0,1)$ ont la même aire $q_A$ et les coûts
modulaires $90\Delta$ et $120\Delta$. Leur différence provient de
l’incidence sectorielle que conserve l’état, outre son occupation totale.

La purification de \ref{prop:ii-area-purificacion-producto} conserve
l’état complet et sa bipartition. Dans une réalisation d’horizon de rayon
$R>0$, l’échelle énergétique et l’opérateur de réponse sont
\begin{equation}
 \tau_R=\frac{\hbar_ac}{2\pi R},\qquad
 H_R=\tau_R\Delta D+e_RI,
 \label{eq:ii-ter27-response}
\end{equation}
avec l’origine $e_R$ enregistrée. Si $\delta E_R=\Tr[(\sigma-\rho_A)H_R]$
et $\delta s=s(\sigma)-s(\rho_A)$, l’origine s’annule dans l’accroissement.

\begin{theorem}[Bilan fini et conversion entre sections thermiques]
\label{thm:ii-ter27-horizon}
Pour tout état $\sigma$ du porteur sectoriel,
\begin{equation}
 \delta E_R-\tau_R\delta s
 =\tau_R D(\sigma\Vert\rho_A)\ge0.
 \label{eq:ii-ter27-free-area}
\end{equation}
Si le même transducteur positif $B$ satisfait
$B(T_H(R))=\tau_R$ et
$B(\Theta_{{\rm clk},a})=\epsilon_a/\ln3$, alors
\begin{equation}
 \frac{T_H(R)}{\Theta_{{\rm clk},a}}
 =\frac{L_\alpha}{R}\frac{\ln3}{2\pi}.
 \label{eq:ii-ter27-clock-horizon}
\end{equation}
\end{theorem}
\begin{proof}
Soustraire les espérances de \eqref{eq:ii-ter27-modular} donne
$\delta s=\Delta\Tr[(\sigma-\rho_A)D]-D(\sigma\Vert\rho_A)$.
Multiplier par $\tau_R$ prouve \eqref{eq:ii-ter27-free-area} ; la
positivité et ses conditions ont déjà été démontrées dans
\ref{thm:ii-area-ley-finita}. Une application linéaire positive entre droites
préserve les quotients positifs. Par conséquent, le quotient des températures est
$\tau_R/(\epsilon_a/\ln3)$ ; \eqref{eq:ii-ter27-energia} l’évalue.
\end{proof}

En $R=L_\alpha$, le facteur de conversion est $\ln3/(2\pi)$.
Les deux sections correspondent à des échelles énergétiques différentes et sont
comparées par un transport exact. Les facteurs de Barbero convertissent
l’occupation en aire ; $\Delta D$ conserve le coût de provenance ; $\tau_R$
l’exprime en énergie. L’identité finie conserve en outre le coût
informationnel de l’état d’essai par rapport à sa référence.

\subsection{Référence spectrale et multiplicités conservées}
\label{subsec:ii-ter27-reference}

La boucle bilatérale précédente produit la raison spectrale $r=1/10$, et le
regroupement tritique conserve $q=r^3=1/1000$ et le reste de chaque indice.
Le premier retour strict au régime de capacité, après la première
vacance, fixe $C_{24}=23$. Le porteur maintient ses trois secteurs de
dimensions $d_0=1$, $d_1=380$, $d_2=649$ ; leurs constructions
d’incidence, de double graduation et de mémoire ont été démontrées dans la section
précédente. Nous conservons l’assignation qui y est déclarée de ces secteurs aux
degrés $23+3n$. Cette assignation est une donnée explicite de la réalisation
conjointe, distincte du calcul qui produit les trois dénombrements.

Dans l’espace $\mathcal D=\bigoplus_{n=0}^2\mathbb C^{d_n}$, soit
\begin{equation}
 \eta=\bigoplus_{n=0}^2r^{23+3n}I_{d_n},\qquad
 b_*:=\Tr\eta
 =10^{-23}\left(1+\frac{380}{1000}+\frac{649}{1000^2}\right).
 \label{eq:ii-ter27-eta}
\end{equation}
L’évaluation exacte donne $b_*=1.380649\times10^{-23}$ et $0<b_*<1$.
L’opérateur conserve chaque étiquette du porteur ; la trace réunit leurs poids.
Pendant les variations thermodynamiques suivantes, $\eta$ reste fixe :
c’est la référence construite, tandis que l’état du système d’essai varie.
Les opérations, les degrés et la normalisation sont ainsi fixés avec
la référence qui intervient dans la variation.

\subsection{Accroissement entropique du registre tensoriel}
\label{subsec:ii-ter27-tensor}

Soit $\mathcal H$ un espace de dimension finie et $\rho$ une matrice de
densité sur celui-ci. Nous écrivons $s(\rho)=-\Tr(\rho\ln\rho)$ et, pour
tout opérateur positif $X$ de trace finie,
$\mathscr H(X)=-\Tr(X\ln X)$, avec $0\ln0=0$. La seconde fonction
admet des références dont la trace diffère de un. Le registre d’essai
$\eta\otimes\rho$ conserve séparément la référence et le système.

\begin{theorem}[Entropie marquée du registre produit]
\label{thm:ii-ter27-tensor}
On a exactement
\begin{equation}
 \mathscr H(\eta\otimes\rho)-\mathscr H(\eta)=b_*s(\rho).
 \label{eq:ii-ter27-tensor}
\end{equation}
\end{theorem}
\begin{proof}
Soient $a_i$ et $p_j$ les valeurs propres respectives. Celles du produit sont
$a_ip_j$, de sorte que
\[
 -\sum_{i,j}a_ip_j\ln(a_ip_j)
 =-\sum_i a_i\ln a_i\sum_jp_j
   -\sum_i a_i\sum_jp_j\ln p_j.
\]
Les identités $\sum_jp_j=1$ et $\sum_i a_i=b_*$ prouvent la formule.
La convention $0\ln0=0$ inclut les états de rang incomplet par continuité.
\end{proof}

La différence sépare la contribution informationnelle de la référence dans
un bilan enregistré. L’opération tensorielle spécifie l’indépendance
de l’essai par rapport à cette référence ; un état corrélé nécessite
l’information conditionnelle correspondante. Le coefficient $b_*$ apparaît
par factorisation de la trace du registre construit.

\subsection{Réalisation normalisée et lecture énergétique conditionnée}
\label{subsec:ii-ter27-flag}

La référence admet une extension normalisée dans
$\widehat{\mathcal D}=\mathcal D\oplus\mathbb C f$ :
\begin{equation}
 \widehat\eta=\eta\oplus(1-b_*)|f\rangle\langle f|,
 \qquad \Omega_\rho=\widehat\eta\otimes\rho.
 \label{eq:ii-ter27-flag}
\end{equation}
Soit $P$ le projecteur sur le facteur direct $\mathcal D$, étendu par l’identité
de $\mathcal H$, et soit $H=H^*$ l’opérateur énergétique du système,
avec origine enregistrée. La marque $f$ complète la trace de cette réalisation
auxiliaire. L’archive TPK reçue, avec ses chemins et sa mémoire, conserve son
propre facteur ; elle est transportée avec le registre et n’est pas remplacée par $f$.

\begin{proposition}[Deux fonctionnelles d’un état normalisé]
\label{prop:ii-ter27-readers}
Avec les logarithmes restreints au support, les lectures
\begin{align}
 \mathcal S_P(\rho)&=-\Tr[P\Omega_\rho(I\otimes\ln\rho)]
                    =b_*s(\rho),\label{eq:ii-ter27-S}\\
 \mathcal E_P(\rho)&=
 \frac{\Tr[P\Omega_\rho(I\otimes H)]}{\Tr(P\Omega_\rho)}
                    =\Tr(\rho H)\label{eq:ii-ter27-E}
\end{align}
préservent respectivement l’information relative à la marque et l’énergie
par état conditionné à celle-ci. Le dénominateur est $\Tr(P\Omega_\rho)=b_*$.
\end{proposition}
\begin{proof}
Sur le facteur direct marqué, $P\Omega_\rho=\eta\otimes\rho$.
La factorisation de la trace donne $b_*s(\rho)$ et
$b_*\Tr(\rho H)$ dans les deux numérateurs. Diviser le second par
$b_*>0$ prouve la lecture énergétique. Tant $\eta$ que $\rho$ se
retrouvent à partir de l’état étendu et du marqueur : la première par sa réduction
marquée et la seconde par la trace sur $\widehat{\mathcal D}$.
\end{proof}

Le bloc complémentaire contient également $\rho$. La lecture informationnelle
totale satisfait $s(\Omega_\rho)-s(\widehat\eta)=s(\rho)$ et son
énergie moyenne totale est $\Tr(\rho H)$. Par conséquent, choisir la fonctionnelle
marquée et l’énergie conditionnée spécifie deux observations d’un état
qui conserve les deux blocs. Le choix des observables et l’élimination
physique d’information sont des opérations différentes.

\subsection{Transducteur fixé par la règle conjointe de lecture}
\label{subsec:ii-ter27-transducer}

Nous travaillons dans les bases positives transportées des droites d’énergie,
de température et d’entropie, avec $u_S=u_E/u_\Theta$. Dans les formules de
coordonnées, $H$ exprime l’énergie en $u_E$ et $\mathcal S_P$ exprime
l’entropie en $u_S$. La règle de réalisation comporte quatre opérations :
recevoir la référence fixe \eqref{eq:ii-ter27-eta} ; évaluer
\eqref{eq:ii-ter27-S} ; évaluer \eqref{eq:ii-ter27-E} ; et définir la
section thermique comme conjuguée de ces deux fonctionnelles. Les bases se
transportent avec les opérateurs, sans introduire de température cible.

\begin{theorem}[Sélection thermique dans la réalisation marquée]
\label{thm:ii-ter27-transducer}
Soit $H$ non scalaire et soit
$\rho_\beta=e^{-\beta H}/Z(\beta)$ pour $\beta>0$. La règle précédente
produit
\begin{equation}
 \frac{d\mathcal S_P}{d\mathcal E_P}=b_*\beta,
 \qquad \Theta_P=\frac1{b_*\beta},\qquad
 B_*(\Theta_Pu_\Theta)=\beta^{-1}u_E,
 \label{eq:ii-ter27-transducer}
\end{equation}
où $B_*(u_\Theta)=b_*u_E$ est une unique application linéaire positive.
\end{theorem}
\begin{proof}
Avec $E=\Tr(\rho_\beta H)$, la différentiation de la somme spectrale
finie donne $E'=-\operatorname{Var}_{\rho_\beta}(H)<0$ et
$(\ln Z)'=-E$. De $s(\rho_\beta)=\beta E+\ln Z$, on déduit
$s'=\beta E'$. La référence reste fixe, donc
$\mathcal S_P'=b_*s'$ et $\mathcal E_P'=E'$. Leur quotient et son
inverse donnent les deux premières identités. L’image d’une base positive
détermine une unique application linéaire ; substituer $\Theta_P$ prouve la dernière.
\end{proof}

Le facteur dépend de la paire concrète de lectures. Si l’on utilise
simultanément l’énergie et l’entropie marquées, $(b_*E,b_*s)$, ou toutes deux
conditionnées, $(E,s)$, la dérivée entropique par rapport à l’énergie est
$\beta$. Dans la paire \eqref{eq:ii-ter27-S}--\eqref{eq:ii-ter27-E},
elle est $b_*\beta$. Ces trois identités se vérifient par la même dérivation ;
elles expliquent quelle normalisation détermine le coefficient. Une fois la
règle marquée fixée, multiplier uniquement son canal entropique par un autre facteur change
la règle et cesse de préserver ses hypothèses.

\begin{corollary}[Principe variationnel marqué]
\label{cor:ii-ter27-variational}
Pour $\Theta>0$, la fonctionnelle
\[
 \Phi_\Theta(\rho)=\mathcal E_P(\rho)-\Theta\mathcal S_P(\rho)
 =\Tr(\rho H)-b_*\Theta s(\rho)
\]
possède un unique minimum sur les matrices de densité :
\begin{equation}
 \rho_\Theta^*=
 \frac{\exp[-H/(b_*\Theta)]}{\Tr\exp[-H/(b_*\Theta)]}.
 \label{eq:ii-ter27-minimum}
\end{equation}
\end{corollary}
\begin{proof}
Substituer $\ln\rho_\Theta^*=-H/(b_*\Theta)-\ln Z$ donne
\[
 \Phi_\Theta(\rho)-\Phi_\Theta(\rho_\Theta^*)
 =b_*\Theta D(\rho\Vert\rho_\Theta^*)\ge0.
\]
La référence est de rang plein. L’égalité dans la positivité de
l’entropie relative se produit exactement pour $\rho=\rho_\Theta^*$,
ce qui prouve l’existence et l’unicité du minimum.
\end{proof}

La référence fixe et la paire de fonctionnelles
\eqref{eq:ii-ter27-S}--\eqref{eq:ii-ter27-E} sont les conditions
constitutives de cette réalisation. Sous ces conditions, la trace $b_*$ détermine
le quotient thermométrique et l’état d’équilibre ; l’identité
d’entropie relative démontre l’existence et l’unicité du minimum.

\subsection{Composition avec la section chronologique et ses transports}
\label{subsec:ii-ter27-map}

Considérons la réalisation chronologique de capacité sur
$\mathcal D=\bigoplus_{n=0}^2\mathbb C^{d_n}$, avec
$(d_0,d_1,d_2)=(1,380,649)$ et $N|_{\mathbb C^{d_n}}=nI$.
Nous conservons l’assignation des degrés $C=23I+3N$ de la section
\ref{sec:ii-ter27-equivalencia}. L’opérateur de
\eqref{eq:ii-ter27-hcap} est
\begin{equation}
 H_{{\rm clk},a}:=H_{\rm cap}
 =\epsilon_a\frac{\ln10}{\ln3}(23I+3N),\qquad
 \beta_{{\rm clk},a}=\frac{\ln3}{\epsilon_a}.
 \label{eq:ii-ter27-clock-operator}
\end{equation}
Il est autoadjoint et positif dans cet espace fini. Ses trois valeurs propres
distinctes, proportionnelles à $23,26,29$, prouvent qu’il est non scalaire.
La multiplication et le calcul spectral donnent exactement
\begin{align}
 \beta_{{\rm clk},a}H_{{\rm clk},a}
   &=(23I+3N)\ln10,\nonumber\\
 e^{-\beta_{{\rm clk},a}H_{{\rm clk},a}}
   &=10^{-23}q^N=\eta,\qquad q=10^{-3},\nonumber\\
 \rho_{{\rm clk},a}
   &=\frac{q^N}{Z_N(q)}=\frac{\eta}{b_*},\qquad
 Z_N(q)=1+380q+649q^2.
 \label{eq:ii-ter27-clock-state}
\end{align}
Le facteur d’origine $10^{-23}$ est conservé dans la trace $b_*$ et s’annule
lorsqu’on normalise l’état. Pour appliquer le théorème
\ref{thm:ii-ter27-transducer}, on fixe l’orientation et $\epsilon_a$,
et donc $H_{\rm cap}$ ; la dérivée est prise par rapport à $\beta$ et
est ensuite évaluée en $\beta_{{\rm clk},a}$. La référence $\eta$,
indépendante de $\epsilon_a$, reste fixe dans le premier facteur,
tandis que $\rho_\beta$ varie dans une seconde copie de $\mathcal D$ avec
le hamiltonien fixe \eqref{eq:ii-ter27-clock-operator}. Son évaluation en
$\beta=\beta_{{\rm clk},a}$ fournit l’état
\eqref{eq:ii-ter27-clock-state} et l’énergie par nat
$\tau_{{\rm clk},a}=\epsilon_a/\ln3$.

La section relative de Williamson maintient
$H_{W,\mathrm{rel}}=\epsilon_a C$ et
$\tau_W=\epsilon_a/\ln10$. Par conséquent,
$H_{{\rm clk},a}/\tau_{{\rm clk},a}
 =H_{W,\mathrm{rel}}/\tau_W=C\ln10$ :
les exposants et l’état coïncident, et les échelles sont reliées par
$\tau_{{\rm clk},a}/\tau_W=\ln10/\ln3$.
La règle marquée, avec ses deux fonctionnelles fixées, produit
\begin{equation}
 \Theta_{{\rm clk},P,a}=\frac{\epsilon_a}{b_*\ln3},\qquad
 B_*(\Theta_{{\rm clk},P,a}u_\Theta)
 =\frac{h_a}{108t_0\ln3}u_E.
 \label{eq:ii-ter27-section-test}
\end{equation}
La température s’obtient ici à partir de la dérivée entropique prouvée, et non
en utilisant le résultat numérique comme condition d’ajustement.

Soient $J_E$ le transport effectif d’énergie vers ces coordonnées et
$J_\Theta$ le transport entre les sections thermométriques du même
état de référence. En particulier, $J_\Theta$ envoie la section
chronologique précédente sur la section de \eqref{eq:ii-ter27-section-test},
selon ses deux lecteurs de température. La preuve d’unicité entre droites
fournit le carré
\begin{equation}
 B_*J_\Theta=J_EB_{\rm clk}.
 \label{eq:ii-ter27-intertwiner}
\end{equation}
En effet, les deux applications coïncident sur la section positive de
l’horloge par \eqref{eq:ii-ter27-section-test}, et cette section engendre la droite.
Le transport déclare les deux réalisations comparées : l’égalité
typée conserve leurs bases et leurs sections ; elle n’identifie pas par simple
notation deux coordonnées thermiques fixées de manière différente.

Pour l’horizon, on obtient
$\Theta_{H,P}(R)=\tau_R/b_*$. Substituer dans
\eqref{eq:ii-ter27-clock-horizon} maintient exactement le quotient entre
les deux températures. De plus, avec $S_P=b_*s$, l’identité
\eqref{eq:ii-ter27-free-area} prend la forme
\[
 \delta E_R-\Theta_{H,P}(R)\,\delta S_P
 =\tau_RD(\sigma\Vert\rho_A).
\]
Cette composition incorpore le nouveau lecteur à un bilan préalable sans
altérer l’incidence, l’action, l’aire ni le terme d’entropie relative.

\subsection{Origine énergétique, réciprocité et section électronique}
\label{subsec:ii-ter27-electron}

La référence énergétique intervient également dans les applications radiales.
Pour un caractère positif $Y$ transporté, on conserve
\[
 H_a=\epsilon_aY,\quad M_a=H_a/c^2,\quad
 R_a=L_\alpha Y,\quad \Lambda_a=L_\alpha Y^{-1}.
\]
Ainsi, $H_a\Lambda_a=\hbar_acI$, $R_a\Lambda_a=L_\alpha^2I$ et
$R_a=(G_a/c^2)M_a$. Un décalage $Y\mapsto Y+\delta I$,
avec positivité conservée, change la masse et le rayon respectivement de
$\epsilon_a\delta/c^2$ et $L_\alpha\delta$, bien que le décalage
scalaire du hamiltonien s’annule dans son état thermique normalisé.
La provenance marquée de l’énergie distingue ces réalisations.

Dans le registre de capacité, l’origine reçue est
$Y_{\rm cap}-Y_{\rm rel}=23\ln10/\ln3$.
La différence énergétique est donc
$\epsilon_a(23\ln10/\ln3)I$ ; sa contribution radiale est conservée
dans les formules précédentes. Si une distribution est reconstruite à partir de
$-\ln\rho$, il faut également restituer $\ln Z$ avant de lire son énergie
absolue. Le marqueur permet de conserver les deux données.

Le correcteur électronique déjà construit vérifie
\begin{align}
 \kappa_e&=80-54\alpha+6\Delta_4,&
 \Delta_4&=\frac{\pi-e\ln\pi}{270}
                 \left(1+\frac\pi{729}\right),\nonumber\\
 \frac{\mathcal E_e^{\beta}}{\mathcal E_e^0}
 &=R_{\rm act}^{\kappa_e},&
 R_{\rm act}&=\frac{\hbar_{\rm pre}}{\hbar_{\rm ret}}.
 \label{eq:ii-ter27-electron-return}
\end{align}
Les énergies électroniques sont exprimées dans la même unité et
$\kappa_e>0$ dans la section électronique construite. Le quotient
positif et son logarithme fixent la puissance réelle. Avec $L_\alpha$, $c$ et
le transducteur communs, les identités précédentes impliquent
\begin{equation}
 \left(\frac{\mathcal E_e^{\beta}}{\mathcal E_e^0}\right)^{1/\kappa_e}
 =R_{\rm act}
 =\frac{\Theta_{{\rm clk,pre}}}{\Theta_{{\rm clk,ret}}}
 =\frac{G_{\rm ret}}{G_{\rm pre}}.
 \label{eq:ii-ter27-return-ratios}
\end{equation}
Pour la branche de retour, soient $\mathcal E_e=m_ec^2$ et
$y_e=\mathcal E_e/\epsilon_{\rm ret}$. Une substitution directe produit
\begin{equation}
 \Lambda_e=\frac{\hbar_{\rm ret}c}{\mathcal E_e},\qquad
 \frac{G_{\rm ret}m_e^2}{\hbar_{\rm ret}c}=y_e^2,\qquad
 \beta_{{\rm clk,ret}}\mathcal E_e=(\ln3)y_e.
 \label{eq:ii-ter27-electron-gravity}
\end{equation}
En effet, $\epsilon_{\rm ret}=\hbar_{\rm ret}c/L_\alpha$ et
$G_{\rm ret}=c^3L_\alpha^2/\hbar_{\rm ret}$ annulent les échelles
dans le second quotient, tandis que le troisième utilise la définition de
$\beta_{{\rm clk,ret}}$. Ces relations sont des compositions du même
retour, avec la conversion dimensionnelle électronique conservée.

L’assignation des degrés $C=23I+3N$, la référence spectrale fixe $\eta$
et la paire de fonctionnelles marquée--conditionnée déterminent conjointement
la réalisation thermique étudiée. Sa preuve tensorielle et variationnelle
établit le coefficient et l’équilibre ; la composition chronologique,
aréale et radiale conserve l’état, les échelles et l’origine énergétique
nécessaires aux lectures respectives.
